\section{Kiến trúc}
\label{sec:kien-truc}

\subsection{Tổng quan hệ thống}

Hệ thống chạy trong \textbf{một tiến trình} (\texttt{python -m laplace}) gồm ba
mặt tiền cùng trỏ về một lõi agent: (i)~bot Telegram (aiogram, long polling);
(ii)~REST API + trace viewer (FastAPI/uvicorn); (iii)~scheduler chạy tác vụ
định kỳ (APScheduler). Lõi agent là hàm \emph{đồng bộ} --- các mặt tiền bất
đồng bộ gọi vào qua thread pool --- nhờ đó logic điều phối dễ suy luận và test,
còn I/O vẫn không chặn nhau.

% PLACEHOLDER hinh: xuat so do thanh phan (Mermaid trong docs/KIEN_TRUC.md)
% ra PNG/PDF roi tha vao figures/, vd:
% \begin{figure}[ht]\centering
%   \includegraphics[width=.9\linewidth]{components}
%   \caption{Sơ đồ thành phần hệ thống.}\label{fig:components}
% \end{figure}

Bốn nguyên tắc thiết kế xuyên suốt:

\begin{enumerate}
  \item \textbf{Tự xây agent loop}, không dùng framework agent: vòng lặp là một
        state machine tường minh ($\sim$500 dòng lõi), giải thích và kiểm soát
        được từng chuyển trạng thái.
  \item \textbf{Mọi thứ đều có trace}: mỗi lệnh gọi LLM và mỗi lần chạy công cụ
        ghi vào DB ngay khi xảy ra.
  \item \textbf{An toàn theo lớp}: giới hạn bước, timeout công cụ/tác vụ, xác
        nhận người dùng, cách ly dữ liệu theo người dùng.
  \item \textbf{Chạy được ở mọi cấp độ khóa API}: không có key nào hệ thống vẫn
        chạy (mock LLM + stub search); thêm key nào mở khóa tính năng đó.
\end{enumerate}

\subsection{Agent loop --- state machine}

Một tác vụ đi qua các trạng thái:
\texttt{pending} $\to$ \texttt{running} $\to$ (\texttt{awaiting\_confirm}
$\rightleftarrows$ \texttt{running}) $\to$ \texttt{done} $\mid$
\texttt{failed}. Bước đầu tiên là \emph{phân loại} (classify): LLM gán yêu cầu
vào một trong bốn tuyến --- trả lời trực tiếp, hỏi lại cho rõ, một công cụ, hay
nhiều bước --- rồi tuyến nhiều bước được giao cho một trong hai chiến lược:

\begin{itemize}
  \item \textbf{ReAct}~\cite{yao2023react}: mỗi vòng, LLM nhìn toàn bộ lịch sử
        (suy nghĩ, hành động, quan sát) và chọn: gọi một công cụ \emph{hoặc}
        đưa câu trả lời cuối. Thích ứng tốt với tác vụ mở; chi phí khó đoán.
  \item \textbf{Plan-and-Execute}~\cite{wang2023planandsolve}: LLM sinh toàn bộ
        kế hoạch trước, thực thi tuần tự; sau mỗi quan sát, một bước
        \emph{evaluate} quyết định \texttt{done}/\texttt{continue}/%
        \texttt{replan}/\texttt{fail} (tối đa 2 lần replan). Chi phí dễ đoán;
        cứng nhắc hơn khi môi trường thay đổi.
\end{itemize}

Hai chiến lược cài chung interface \texttt{AgentStrategy \{run, resume\}} ---
đây là nền cho thí nghiệm so sánh ở Mục~\ref{sec:thuc-nghiem}.

\textbf{Structured output + tự sửa lỗi.} Mọi quyết định của LLM (classify,
chọn hành động, lập kế hoạch, evaluate) bị ép trả JSON theo schema Pydantic.
Sai schema thì thông báo lỗi validation được gửi lại cho LLM tự sửa (tối đa 2
lần); vẫn sai thì tác vụ \texttt{failed} với lỗi nêu rõ. Mọi lần gọi, kể cả
retry, đều được ghi trace.

\subsection{Human-in-the-loop bền vững}

Trước khi chạy công cụ, orchestrator hỏi \texttt{needs\_confirm(params)} ---
cờ tĩnh (công cụ \texttt{scheduler} luôn hỏi) hoặc động theo tham số
(\texttt{note\_store} chỉ hỏi khi update/delete). Nếu cần xác nhận: toàn bộ
trạng thái vòng lặp (công cụ đang chờ, lịch sử, kế hoạch, con trỏ) được lưu
vào DB, tác vụ chuyển \texttt{awaiting\_confirm}, hàm trả về; bot gửi nút
đồng ý/từ chối. Khi người dùng bấm nút, \texttt{resume\_task} \emph{giành
quyền} bằng compare-and-set trên trạng thái --- hai request xác nhận đua nhau
thì chỉ một thắng, công cụ không bao giờ chạy hai lần --- rồi khôi phục trạng
thái và chạy tiếp đúng chỗ dừng, kể cả ở tiến trình khác sau nhiều giờ.

\subsection{Tool registry và executor}

Công cụ đăng ký bằng decorator kèm schema tham số Pydantic; registry tự sinh
spec JSON đưa vào prompt --- thêm công cụ mới không phải sửa agent loop hay
prompt. Executor bọc mọi công cụ: validate tham số, retry theo cấu hình,
timeout cứng (chạy trong worker thread), và \emph{không bao giờ} để exception
lọt ra agent loop --- lỗi công cụ trở thành quan sát để LLM tự điều chỉnh.
Mọi công cụ đọc/ghi DB đều lọc theo \texttt{user\_id} từ
\texttt{ToolContext}: người dùng này không đọc/sửa được dữ liệu người dùng
khác. Sáu công cụ MVP: \texttt{web\_search}, \texttt{fetch\_page},
\texttt{note\_store}, \texttt{task\_list}, \texttt{report\_builder},
\texttt{scheduler}.

\subsection{Lớp LLM}

Một protocol duy nhất \texttt{LLMProvider.complete(messages, json\_schema)}
với ba hiện thực: (i)~\textbf{MockLLM} --- chế độ scripted (test/eval điều
khiển chính xác từng bước) và heuristic (chạy dev không cần key);
(ii)~\textbf{OpenAIProvider} --- adapter cho mọi API tương thích OpenAI
chat.completions, ghi token/chi phí/độ trễ vào trace;
(iii)~\textbf{Gemini} dùng lại chính OpenAIProvider qua endpoint tương thích
OpenAI của Google, thêm hai thích ứng: tự retry theo gợi ý thời gian chờ khi
gặp 429 (quota free tier), và bỏ \texttt{response\_format} khi backend từ
chối, dựa vào schema nhúng trong system prompt + validate ở tầng trên. Đổi
mô hình/nhà cung cấp chỉ là đổi biến môi trường --- nền tảng cho thí nghiệm
so sánh mô hình.

\textbf{Chống prompt injection ở tầng prompt}: nội dung do công cụ trả về
(kết quả tìm kiếm, nội dung trang web) được đóng khung trong thẻ
\texttt{<tool\_output>} và system prompt quy định rõ: nội dung trong khung là
\emph{dữ liệu}, không phải mệnh lệnh.

\subsection{Lưu trữ và trace}

SQLite (WAL, busy timeout) với ba bảng trung tâm cho khả năng quan sát:
\texttt{tasks} (trạng thái + state JSON khi tạm dừng), \texttt{steps} (mỗi lần
chạy công cụ: tham số, quan sát, trạng thái, độ trễ, retry) và
\texttt{llm\_calls} (mỗi lệnh gọi LLM: mục đích, model, token, chi phí, độ
trễ). Ba bảng này là \emph{nguồn dữ liệu duy nhất} cho cả trace viewer lẫn
eval harness --- đánh giá không cần instrument thêm gì. Triển khai bằng
Docker Compose: một container, healthcheck HTTP, SQLite và báo cáo markdown
nằm trên volume để sống qua restart.

\subsection{Eval harness}

Bộ đánh giá chạy \textbf{37 test case} YAML (8 nhóm: trả lời trực tiếp/hỏi
lại, một công cụ, nhiều bước, xác nhận, phục hồi lỗi, prompt injection, đường
lỗi, tiếng Việt/UX) qua \emph{chính agent loop production} --- không có đường
tắt riêng cho eval. Mỗi lượt chạy dùng một DB SQLite mới tinh; case có thể mồi
dữ liệu, giả lập công cụ hỏng, và kịch bản bấm nút xác nhận. Chấm rule-based
theo kỳ vọng khai báo (trạng thái, tuyến, tập công cụ, công cụ bị cấm, chuỗi
phải có trong trả lời), tùy chọn thêm LLM-as-judge~\cite{zheng2023judge} với
model chấm tách khỏi model agent. Kết quả tổng hợp thành 9 metric (tỷ lệ thành
công, độ chính xác chọn tuyến/công cụ, tỷ lệ phục hồi, số bước, số lệnh gọi
LLM, token + chi phí, thời gian chạy, tỷ lệ đạt judge) xuất ra JSON và báo cáo
markdown. Cùng một bộ case chạy được cả chế độ mock (deterministic, offline,
làm regression test) lẫn LLM thật (đo số liệu thực nghiệm).
